\scanpage{78}
\mspage{77}

En effet,
\[
\unsure{C_M|c_M}(A)=\{X\in M\mid X\parallel A\}
=\{X\in M\mid |X|\cap|A|=\varnothing
\]
\inserted{i.e.\ $|\overline X|\subset\unsure{R|P}\setminus A$}
donc
\[
|C_M(A)|=\bigcup\overline X\quad\text{pour }X\unsure{\in|\notin}\deleted{A}M\ \unsure{\text{t.q.}|\text{i.e.}}\quad |X|\subset P\setminus A
\]
d'où l'inclusion dite, sans garantie que ce soit tjrs égalité.

On aimerait aussi avoir, pour $A$ partie quelconque de $M$
\[
|A|\overset{?}{=}|\operatorname{supp}A|,
\tag*{(**)}
\]
alors qu'on a a priori (tautologiquement)
\[
|A|\subset|\operatorname{supp}A|
\]
puisque $A\subset\operatorname{supp}A$. \deleted{\unsure{le}} Donc on voudrait
\[
|\operatorname{supp}A|\subset|A|
\]
ce qui serait conséquence de l'implication inverse (hypothétique) de celle du Lemme. La validité de cette implication inverse, pour $\forall A,B$, est d'ailleurs équivalente au cas particulier d'un couple $(\operatorname{supp}A,A)$, i.e.\ à la validité de (**).

Cela fait donc quelques \unsure{perplexités}, \unsure{liées|liés} aux \unsure{mystères} des \deleted{de} $P$-supports.

Tout \inserted{\unsure{ça}} se \unsure{résoud|résout} si on fait l'hypothèse suivante

\uline{Hyp} Soit $A\subset M$, et $x\in P\setminus|A|$. Alors $\exists X\in M$ avec \deleted{$\unsure{X}$} $X\parallel A$ (i.e.\ $|X|\subset P\setminus|A|$) et $x\in|X|$
